\documentclass[11pt,a4paper]{article}

% --- 核心宏包与中文支持 ---
\usepackage[utf8]{inputenc}
\usepackage{xeCJK}
\usepackage{amsmath, amssymb, amsfonts}
\usepackage{listings}
\usepackage{xcolor}
\usepackage{geometry}
\usepackage{tcolorbox}
\usepackage{booktabs}
\usepackage{tabularx}
\usepackage{setspace}
\usepackage{enumitem}
\usepackage[hidelinks]{hyperref}
\usepackage{bookmark}
\usepackage{fancyhdr}

% --- PDF 书签与元数据 ---
\hypersetup{
    colorlinks=false,
    pdfborder={0 0 0},
    bookmarks=true,
    bookmarksnumbered=true,
    bookmarksopen=true,
    bookmarksopenlevel=2,
    pdfcreator={XeLaTeX},
    pdftitle={09 - Pandas Library II},
    pdfauthor={黄子扬},
}

% --- 页面美化设置 ---
\geometry{top=0.7in,bottom=0.8in,left=0.6in,right=0.6in}
\onehalfspacing
\tcbuselibrary{skins,breakable}
\setlength{\parindent}{0pt}
\setlength{\parskip}{0.6em}

% --- 代码样式配置 ---
\definecolor{mainblue}{RGB}{0,51,102}
\definecolor{examred}{RGB}{180,0,0}
\definecolor{codebg}{RGB}{248,248,248}
\definecolor{commentgreen}{RGB}{0,128,0}

\lstset{
    backgroundcolor=\color{codebg},
    basicstyle=\ttfamily\footnotesize,
    keywordstyle=\color{blue}\bfseries,
    commentstyle=\color{commentgreen},
    stringstyle=\color{orange},
    breaklines=true,
    showstringspaces=false,
    frame=leftline,
    xleftmargin=2em,
    numbers=left,
    numbersep=5pt,
    numberstyle=\tiny\color{gray},
    language=Python,
    morekeywords={None, True, False, pd, np, lambda, inplace, astype, cut, groupby, pivot_table, assign, import, as, for, in, if}
}

% --- 自定义教学框 ---
\newtcolorbox{concept}[1]{colback=blue!5, colframe=mainblue, fonttitle=\bfseries, title={{【Concept / 核心概念】#1}}, arc=3pt, breakable}
\newtcolorbox{warning}[1]{colback=red!5, colframe=examred, fonttitle=\bfseries, title={{【Warning / 避坑指南】#1}}, arc=3pt, breakable}
\newtcolorbox{examplebox}[1]{colback=green!2, colframe=green!40!black, fonttitle=\bfseries, title={{【Example / 深度案例：house.csv】#1}}, arc=2pt, breakable}

% --- 页眉页脚 ---
\pagestyle{fancy}
\fancyhf{}
\fancyhead[L]{\small\bfseries 09 - Pandas Library II / Pandas 库 II}
\fancyhead[R]{\small\thepage}
\renewcommand{\headrulewidth}{0.5pt}

% --- 文档信息 ---
\title{\Huge \bfseries 09 - Pandas Library II \\ Data Exploration \& Feature Engineering}
\author{\Large \textbf{哈尔滨工业大学 \quad 黄子扬}}
\date{2026年6月8日}

\begin{document}

\maketitle
\tableofcontents
\newpage

% ======================================================================
\section{Setup \& Data Loading / 环境配置与数据载入}
% ======================================================================

\begin{examplebox}{Initial Setup / 初始化配置}
\begin{lstlisting}
import pandas as pd
import numpy as np
import matplotlib.pyplot as plt

# Pandas display options / 显示选项
pd.set_option('display.max_columns', 15)
pd.set_option('display.max_rows', 15)
pd.set_option('display.float_format', lambda x: '%.3f' % x)
\end{lstlisting}
\end{examplebox}

\subsection{Loading Data / 数据载入}

\begin{concept}{Universal Readers / 万能读取器}
    Pandas supports various formats:
    \begin{itemize}[itemsep=0.3em]
        \item \textbf{\texttt{pd.read\_csv()}}：Most common — loads comma-separated data.
        \item \textbf{\texttt{pd.read\_excel()}}：Loads data from Excel sheets.
        \item \textbf{\texttt{pd.read\_html()}}：Scrapes tables directly from a website URL.
        \item \textbf{\texttt{pd.read\_json()}}：Read from JSON formatted string, URL or file.
        \item \textbf{\texttt{pd.read\_table()}}：From a delimited text file (like TSV).
    \end{itemize}
\end{concept}

\begin{examplebox}{Reading the house.csv Dataset / 读取房屋数据集}
\begin{lstlisting}
data = pd.read_csv('data/house.csv')
print(data.shape)       # (rows, columns) — check dataset size
data.head(5)            # View first 5 rows
data.tail(5)            # View last 5 rows
\end{lstlisting}
\end{examplebox}

% ======================================================================
\section{Data Cleaning / 数据清洗（实战灵魂）}
% ======================================================================

\subsection{Dropping Redundant Columns / 删除冗余列}

\begin{concept}{Why Drop Columns? / 为什么要删列？}
    Columns that have the same value for all rows (like 'country' or 'postcode') provide \textbf{no information} — they should be dropped. \\
    对所有行都相同的列（如国家、邮编）不提供任何有效信息，应该删除。
\end{concept}

\begin{examplebox}{Dropping Columns / 删除列}
\begin{lstlisting}
print(data.columns)                     # See all column names

# Drop columns that won't be used
data = data.drop(columns=["country", "postcode"])
print(data.shape)                       # Reduced column count
data.head(5)
\end{lstlisting}
\end{examplebox}

\subsection{Handling NaNs / 处理缺失值}

\begin{warning}{NaN Detection and Deletion / 缺失值检测与删除}
    \begin{itemize}[itemsep=0.8em]
        \item \textbf{\texttt{data.isna().sum()}}：Counts missing values in each column. \texttt{isnull()} is equivalent.
        \item \textbf{\texttt{data.dropna(inplace=True)}}：Removes \textbf{any row} containing a missing value. \texttt{inplace=True} modifies the original DataFrame.
        \item \textbf{\texttt{df.fillna(value)}}：Replaces all null values with a specific value (e.g., the mean, or 0).
        \item \textbf{axis variants}：\texttt{dropna(axis=1)} drops columns with nulls; \texttt{dropna(thresh=n)} keeps rows with at least n non-null values.
    \end{itemize}
\end{warning}

\begin{examplebox}{NaN Handling Example / 缺失值处理示例}
\begin{lstlisting}
# Check for NaNs / 检查缺失值
print(data.isna().sum())
print(data.isnull().sum())              # Same as isna()

# Remove rows with any NaN / 删除含 NaN 的行
data.dropna(inplace=True)
print(data.shape)                       # Fewer rows after dropping

# Alternative: fill NaNs with mean / 替代方案：用均值填充
# data['column'].fillna(data['column'].mean(), inplace=True)
\end{lstlisting}
\end{examplebox}

\subsection{Renaming Columns / 重命名列}

\begin{examplebox}{Renaming for Convenience / 简化列名}
    Long column names are hard to type repeatedly. Rename them for convenience.
\begin{lstlisting}
data.rename(columns={
    'rentEstimate_currentPrice': 'rentPrice',
    'saleEstimate_currentPrice': 'salePrice'
}, inplace=True)
data.head(5)
\end{lstlisting}
\end{examplebox}

\subsection{Type Conversion / 类型转换}

\begin{concept}{Optimizing Data Types / 优化数据类型}
    Some data types (like \texttt{object}) should be \texttt{category} for memory efficiency. \\
    将只有几个不同值的字符串列转为 \texttt{category} 类型可以大幅节省内存。
\end{concept}

\begin{examplebox}{Type Conversion Example / 类型转换示例}
\begin{lstlisting}
print(data.dtypes)                      # Check current types

# Convert 'propertyType' from object to category
data['propertyType'] = data.propertyType.astype('category')
data.head(5)
\end{lstlisting}
\end{examplebox}

% ======================================================================
\section{Feature Engineering / 特征工程（进阶技巧）}
% ======================================================================

\subsection{lambda Functions / 匿名函数}

\begin{concept}{What is a lambda? / 什么是 lambda？}
    Lambda functions are small, \textbf{one-line}, \textbf{anonymous} functions. They can receive multiple arguments but can only contain \textbf{one expression} (the return value). \\
    Lambda 是小型的一行匿名函数，只能包含一个表达式。

    \begin{lstlisting}
# Syntax: lambda arguments: expression
square = lambda x: x**2
print(square(5))    # 25
    \end{lstlisting}
\end{concept}

\subsection{Creating New Columns / 创建新列}

\begin{examplebox}{Boolean Flag Column / 布尔标记列}
\begin{lstlisting}
# Create a boolean column 'above_mean'
data['above_mean'] = data.salePrice > data.salePrice.mean()
data.head(5)
\end{lstlisting}
\end{examplebox}

\begin{examplebox}{Using assign() / 使用 assign() 方法}
    \texttt{assign()} also creates new columns using lambda:
\begin{lstlisting}
data.assign(above_mean=lambda x: x.salePrice > x.salePrice.mean())
data.head(5)
\end{lstlisting}
\end{examplebox}

\subsection{Filtering Data / 数据过滤}

\begin{examplebox}{Multi-Condition Filtering / 多条件过滤}
    Use \texttt{\&} (and), \texttt{|} (or) between conditions. Each condition must be wrapped in \textbf{parentheses}!
\begin{lstlisting}
# Filter: flat type, bedrooms >= 2, bathrooms <= 2
filtered = data.loc[(data['propertyType'] == 'Purpose Built Flat') &
                    (data['bedrooms'] >= 2.0) &
                    (data['bathrooms'] <= 2.0)]
print(filtered)
\end{lstlisting}
\end{examplebox}

\begin{warning}{Use \& not 'and' for Element-Wise Logic / Pandas 中要用 \& 而不是 and}
    \texttt{and}/\texttt{or} work on single booleans. \texttt{\&}/\texttt{|} work element-wise on Series/arrays. Every condition MUST be wrapped in parentheses: \texttt{(cond1) \& (cond2)}.
\end{warning}

\subsection{Conditional Modification / 条件修改}

\begin{examplebox}{Modifying Values Based on Condition / 条件修改值}
\begin{lstlisting}
# Check value distribution / 查看值分布
print(data['above_mean'].value_counts())

# Swap True/False labels / 交换标签
data.loc[data['above_mean'] == False, 'above_mean'] = 'yes'
data.loc[data['above_mean'] == True, 'above_mean'] = 'no'
data.head(5)
\end{lstlisting}
\end{examplebox}

\subsection{Sorting by Multiple Columns / 多列排序}

\begin{examplebox}{Multi-Column Sort / 多列排序}
\begin{lstlisting}
# Sort salePrice ascending, bathrooms descending
data.sort_values(['salePrice', 'bathrooms'], ascending=[True, False])
# Equivalent: ascending=[1, 0] where 1=ascending, 0=descending
\end{lstlisting}
\end{examplebox}

\subsection{Discretization (Binning) / 数据离散化（分箱）}

\begin{concept}{pd.cut() / 分箱操作}
    Converts a continuous numerical variable into discrete "bins" (categories). \\
    将连续数值变量划分为离散的"箱子"（分类）。
    \texttt{pd.cut(data, bins=3, labels=['low', 'medium', 'high'])}
\end{concept}

\begin{examplebox}{Price Binning Example / 价格分箱示例}
\begin{lstlisting}
# Divide salePrice into 3 equal-width bins
data["salePriceBins"] = pd.cut(data["salePrice"], bins=3,
                                labels=['low', 'medium', 'high'])
data.head(5)
print(data["salePriceBins"].value_counts())  # Count per bin
\end{lstlisting}
\end{examplebox}

\subsection{Encodings / 类别数据编码}

\begin{concept}{Two Encoding Methods / 两种编码方式}
    \begin{enumerate}[itemsep=0.5em]
        \item \textbf{LabelEncoder}：Converts strings into integers (e.g., 'Freehold' $\rightarrow$ 0, 'Leasehold' $\rightarrow$ 1). Use when there is an \textbf{ordinal} relationship.
        \item \textbf{One Hot Encoding (\texttt{pd.get\_dummies})}：Creates N new binary columns for N categories. Each column contains 0 or 1. Use when categories are \textbf{nominal} (no order).
    \end{enumerate}
\end{concept}

\begin{examplebox}{LabelEncoder Example / 标签编码}
\begin{lstlisting}
from sklearn.preprocessing import LabelEncoder

# Check current values / 查看当前值
print(data.tenure.value_counts())

# Encode: Freehold → 0, Leasehold → 1
tenure_encoder = LabelEncoder()
data['tenure'] = tenure_encoder.fit_transform(data['tenure'])
print(data.head(5))
print(data.tenure.value_counts())

# See the mapping / 查看映射关系
print({k: v for k, v in enumerate(list(tenure_encoder.classes_))})
\end{lstlisting}
\end{examplebox}

\begin{examplebox}{One Hot Encoding Example / 独热编码}
\begin{lstlisting}
# Create dummy columns for 'tenure'
data_encoder = pd.get_dummies(data, columns=["tenure"])
print(data_encoder.shape)           # Now has more columns
data_encoder.head()
\end{lstlisting}
\end{examplebox}

% ======================================================================
\section{Data Exploration \& Aggregation / 数据探索与聚合}
% ======================================================================

\subsection{Group Analysis / 分组分析}

\begin{examplebox}{Exploring by Bathrooms / 按浴室数量分析}
\begin{lstlisting}
print(data["bathrooms"].unique())              # All unique bathroom counts
print(data["bathrooms"].value_counts())         # Frequency of each count

print(data.groupby("bathrooms").size())         # Same as value_counts
print(data.groupby("bathrooms")['salePrice'].mean())  # Avg price per bathroom count
\end{lstlisting}
\end{examplebox}

\subsection{Pivot Table / 透视表}

\begin{concept}{Multi-dimensional Aggregation / 多维汇总}
    A \textbf{pivot table} is a table of grouped values that aggregates individual items into discrete categories. \\
    透视表将数据按照两个维度进行分组并聚合。
\end{concept}

\begin{examplebox}{Pivot Table Example / 透视表示例}
\begin{lstlisting}
# Mean salePrice by bathrooms AND price bins
data.pivot_table(index='bathrooms',
                 columns='salePriceBins',
                 values='salePrice',
                 aggfunc='mean')
\end{lstlisting}
    This creates a 2D table: rows = bathroom count, columns = price bin, values = mean sale price.
\end{examplebox}

\subsection{Crosstabs / 交叉表}

\begin{concept}{Frequency Table / 频率表}
    \texttt{pd.crosstab()} provides an easy way to create a frequency table — counting how many occurrences exist for combinations of categories. \\
    交叉表用于统计两个分类变量的组合出现频率。
\end{concept}

\begin{examplebox}{Crosstab Example / 交叉表示例}
\begin{lstlisting}
pd.crosstab(index=data["bathrooms"], columns=data["salePriceBins"])
\end{lstlisting}
    This shows: for each bathroom count, how many properties fall into each price bin.
\end{examplebox}

\subsection{Describe / 统计摘要}

\begin{examplebox}{General Statistics / 整体统计}
\begin{lstlisting}
data.describe()
# Returns count, mean, std, min, 25%, 50%, 75%, max for all numerical columns
\end{lstlisting}
\end{examplebox}

% ======================================================================
\section{Visualization with Pandas / 数据可视化}
% ======================================================================

\begin{concept}{Integrated Plotting / 集成绘图}
    Pandas provides a \texttt{.plot()} method that directly uses Matplotlib. For SVG-format plots (high resolution), use:
\begin{lstlisting}
%config InlineBackend.figure_formats = ['svg']
%matplotlib inline
\end{lstlisting}
\end{concept}

\subsection{Line Plot / 折线图}

\begin{examplebox}{Line Plot with Pandas / Pandas 折线图}
\begin{lstlisting}
# Plot salePrice against rentPrice for rows 100-110
data.loc[100:110].plot(x="rentPrice", y="salePrice",
                       title='salePrice vs rentPrice',
                       ylabel='salePrice', alpha=0.8)
plt.show()
\end{lstlisting}
\end{examplebox}

\subsection{Bar Plot / 条形图}

\begin{examplebox}{Customized Bar Plot from Pivot Table / 透视表条形图}
\begin{lstlisting}
# Create pivot table for plotting
plot_data = data.pivot_table(index='bathrooms', columns='salePriceBins',
                             values='salePrice', aggfunc='mean')

# Plot as bar chart with custom formatting
ax = plot_data.plot(
    kind='bar', rot=0, xlabel='', ylabel='Price Mean', fontsize=8,
    figsize=(12, 1.5), title='Number of bathrooms | Price Bins',
)
ax.legend(title='', loc='center', bbox_to_anchor=(0.5, -0.3),
          ncol=3, frameon=False)
plt.show()
\end{lstlisting}
\end{examplebox}

% ======================================================================
\section{Summary：Advanced Operators / 进阶操作符速查表}
% ======================================================================

\begin{table}[h]
\centering
\small
\renewcommand{\arraystretch}{1.3}
\begin{tabular}{@{}ll@{}}
\toprule
\textbf{Operator / 操作符} & \textbf{Description / 描述} \\ \midrule
\texttt{df.isna().sum()} & Count missing values per column / 每列空值计数 \\
\texttt{df.isnull().sum()} & Same as isna() / 同上 \\
\texttt{df.dropna(inplace=True)} & Remove rows with any NaN / 删除含 NaN 的行 \\
\texttt{df.fillna(x)} & Replace all null values with x / 用 x 填充空值 \\
\texttt{s.value\_counts()} & View unique values and their counts / 查看各类别出现频率 \\
\texttt{df.sort\_values(c, ascending=False)} & Sort by column in descending order / 按列降序排列 \\
\texttt{pd.merge(df1, df2, how='left')} & SQL-style left join / SQL 风格左连接 \\
\texttt{df.groupby(c).mean()} & Group by column and calculate mean / 分组均值 \\
\texttt{df.pivot\_table(index, columns, values, aggfunc)} & Multi-dimensional aggregation / 透视表多维聚合 \\
\texttt{pd.crosstab(index, columns)} & Frequency table / 交叉频率表 \\
\texttt{pd.cut(data, bins, labels)} & Discretize continuous values / 连续值离散化(分箱) \\
\texttt{pd.get\_dummies(df, columns=[c])} & One-hot encoding / 独热编码 \\
\texttt{df.rename(columns=\{old:new\})} & Rename columns / 列重命名 \\
\texttt{df.astype('category')} & Convert column type for memory efficiency / 类型转换 \\
\texttt{df.assign(col=lambda x: expr)} & Create new column with lambda / 用 lambda 创建新列 \\
\texttt{df.plot(kind='bar')} & Pandas integrated plotting / Pandas 集成绘图 \\
\texttt{\%config InlineBackend.figure\_formats} & Set SVG output for high resolution / SVG 高清显示 \\ \bottomrule
\end{tabular}
\end{table}

\end{document}
